% GENERATED by bin/curate/gen_tutorials_guide.py -- DO NOT EDIT BY HAND.
% Assembled from each case's controlDict description + header comment.

\subsection{ChemicalPlantTutorial}
\grouptheory{ChemicalPlantTutorial}
\clearpage
\tutentry{CONCENTRATION}{tutorials/plant/ChemicalPlantTutorial/CONCENTRATION}
\begin{tutnarrative}
/ system/flowsheetDict (SECTOR -- COMPOSITE)

Double-effect evaporator (vapour of effect 1 heats effect 2) feeding a crystalliser. Same shape as the plant flowsheetDict, one level down: children + boundary + default feeds + connections.

\end{tutnarrative}

\clearpage
\tutentry{DRYING}{tutorials/plant/ChemicalPlantTutorial/DRYING}
\begin{tutnarrative}
/ system/flowsheetDict (SECTOR -- COMPOSITE)

Spray dryer $\to$ additional solid drying $\to$ cyclone dust separator. The solid powder goes SD $\to$ BD $\to$ product; the spray-dryer exhaust air (carrying fines) goes SD $\to$ CY, which returns clean air + recovered dust.

\end{tutnarrative}

\clearpage
\tutentry{FERMENTATION}{tutorials/plant/ChemicalPlantTutorial/FERMENTATION}
\begin{tutnarrative}
/ system/flowsheetDict (SECTOR -- COMPOSITE with INTERNAL RECYCLE)

Continuous fermentation side-loop of the sugar plant: a slip-stream of fresh juice (Must) is hydrolysed + fermented into ethanol + CO2, the Flash splits the volatiles (EthanolVapour product) from the liquid bottoms, and a Splitter recycles 80 \% of the bottoms back to the Mixer so unconverted sucrose gets another shot. Pedagogical purpose:

1. Reaction (Arrhenius first-order in sucrose). 2. Vapour-liquid flash for separation. 3. Tear-stream RECYCLE that the Wegstein/Newton outer loop must iterate to a fixed point.

Must    [Mixer]  MixedFeed   [Fermentor]  FermOut   [Flash]   EthanolVapour           Liquid       Recycle     [Splitter]                     

Tear stream: Recycle (Splitter output $\to$ Mixer input).

\end{tutnarrative}

\clearpage
\tutentry{JuiceSplitter}{tutorials/plant/ChemicalPlantTutorial/JuiceSplitter}
\begin{tutnarrative}
/ system/flowsheetDict (LEAF -- plant-level juice split)

Diverts the fresh sugar-cane juice between the two production lines of the mill: 2/3 goes to the CONCENTRATION line (evaporators $\to$ crystals $\to$ dry sugar) and 1/3 to the FERMENTATION line (yeast $\to$ ethanol). This is what a real Brazilian sugar-ethanol plant does after the clarifier: the same juice header feeds both products, with the split ratio set by market price + crop quality.

Pedagogically the splitter is the single inter-sector cable that ties the two downstream branches together -- without it, FERMENTATION would be a parallel island.

\end{tutnarrative}

\clearpage
\tutentry{acetonePlant}{tutorials/plant/acetonePlant}
\begin{tutscheme}
\begin{tikzpicture}[>={Stealth[length=2mm]}, every node/.style={font=\scriptsize}]
  \node[draw, rounded corners=2pt, align=center, minimum width=1.9cm, minimum height=0.85cm, fill=black!4] (M1) at (237.8,-0.0) {\textbf{M1}\\[-1pt]{\tiny mixer}};
  \node[draw, rounded corners=2pt, align=center, minimum width=1.9cm, minimum height=0.85cm, fill=black!4] (vaporiser) at (240.7,-0.0) {\textbf{vaporiser}\\[-1pt]{\tiny phaseChanger}};
  \node[draw, rounded corners=2pt, align=center, minimum width=1.9cm, minimum height=0.85cm, fill=black!4] (reactor) at (243.6,-0.0) {\textbf{reactor}\\[-1pt]{\tiny conversionReactor}};
  \node[draw, rounded corners=2pt, align=center, minimum width=1.9cm, minimum height=0.85cm, fill=black!4] (separator) at (246.5,-0.0) {\textbf{separator}\\[-1pt]{\tiny isothermalFlash}};
  \node[draw, rounded corners=2pt, align=center, minimum width=1.9cm, minimum height=0.85cm, fill=black!4] (absorber) at (249.4,-0.0) {\textbf{absorber}\\[-1pt]{\tiny absorber}};
  \node[draw, rounded corners=2pt, align=center, minimum width=1.9cm, minimum height=0.85cm, fill=black!4] (M2) at (252.3,-0.0) {\textbf{M2}\\[-1pt]{\tiny mixer}};
  \node[draw, rounded corners=2pt, align=center, minimum width=1.9cm, minimum height=0.85cm, fill=black!4] (vent) at (255.2,-0.0) {\textbf{vent}\\[-1pt]{\tiny isothermalFlash}};
  \node[draw, rounded corners=2pt, align=center, minimum width=1.9cm, minimum height=0.85cm, fill=black!4] (C1) at (258.1,-0.0) {\textbf{C1}\\[-1pt]{\tiny distillationColumn}};
  \node[draw, rounded corners=2pt, align=center, minimum width=1.9cm, minimum height=0.85cm, fill=black!4] (C2) at (261.0,-0.0) {\textbf{C2}\\[-1pt]{\tiny distillationColumn}};
  \draw[->] ($(M1.west)+(-1.05,0.00)$) -- ($(M1.west)+(0,0.00)$) node[at start, above, font=\tiny, anchor=south west] {freshFeed};
  \draw[->, dashed] (C2.south) .. controls +(0,-1.0) and +(0,-1.0) .. (M1.south) node[midway, below, font=\tiny] {ipaRecycle};
  \draw[->] (M1) -- (vaporiser) node[midway, above, font=\tiny] {mixedFeed};
  \draw[->] (vaporiser) -- (reactor) node[midway, above, font=\tiny] {Rin};
  \draw[->] (reactor) -- (separator) node[midway, above, font=\tiny] {Rout};
  \draw[->] (separator) -- (absorber) node[midway, above, font=\tiny] {flashGas};
  \draw[->] ($(absorber.west)+(-1.05,0.00)$) -- ($(absorber.west)+(0,0.00)$) node[at start, above, font=\tiny, anchor=south west] {washWater};
  \draw[->] (separator) -- (M2) node[midway, above, font=\tiny] {flashLiquid};
  \draw[->] (absorber) -- (M2) node[midway, above, font=\tiny] {absBottoms};
  \draw[->] (M2) -- (vent) node[midway, above, font=\tiny] {F1};
  \draw[->] (vent) -- (C1) node[midway, above, font=\tiny] {columnFeed};
  \draw[->] (C1) -- (C2) node[midway, above, font=\tiny] {B1};
  \draw[->] ($(absorber.east)+(0,0.00)$) -- ($(absorber.east)+(1.05,0.00)$) node[at end, above, font=\tiny, anchor=south east] {offgas};
  \draw[->] ($(vent.east)+(0,0.00)$) -- ($(vent.east)+(1.05,0.00)$) node[at end, above, font=\tiny, anchor=south east] {ventGas};
  \draw[->] ($(C1.east)+(0,0.00)$) -- ($(C1.east)+(1.05,0.00)$) node[at end, above, font=\tiny, anchor=south east] {acetoneProduct};
  \draw[->] ($(C2.east)+(0,0.00)$) -- ($(C2.east)+(1.05,0.00)$) node[at end, above, font=\tiny, anchor=south east] {waterProduct};
\end{tikzpicture}
\end{tutscheme}
\tutwhat{Acetone from 2-propanol: Luyben's Figure 1 closed, with the IPA recycle}
\begin{tutnarrative}
ACETONE FROM 2-PROPANOL, THE CLOSED FLOWSHEET

W. L. Luyben, "Design and Control of the Acetone Process via Dehydrogenation of 2-Propanol", Ind. Eng. Chem. Res. 50 (2011) 1206-1218, DOI 10.1021/ie901923a -- his Figure 1 (the Turton base case), end to end:

fresh feed + IPA recycle $\to$ vaporiser $\to$ reactor $\to$ separator $\to$ absorber (offgas washed) $\to$ C1 $\to$ acetone product $\to$ C2 $\to$ water product $\to$ IPA RECYCLE (tear)

This is the case the six preceding ones were built to make possible. Each of them isolates one unit on Luyben's OWN published inlet, so that whatever it gets wrong is that unit and its thermodynamics. Here nothing is published except the fresh feed: every internal stream is the plant's own answer, and the errors compose.

WHAT THIS CASE CLAIMS, AND WHAT IT DOES NOT. It claims to be Luyben's TOPOLOGY and Luyben's EQUIPMENT SPECIFICATIONS solved on the thermodynamics Choupo actually has. It does NOT claim to reproduce his stream table -- the siblings measured, in advance, exactly why it cannot, and the gap between his table and this one is the deliverable rather than a defect to be closed:

acetone01: UNIFAC puts the IPA/water azeotrope at x 0.590 against his 0.673, and the lake-estimated isopropanol boils 4.5 K high. acetone04: UNIFAC INVENTS an acetone/water azeotrope at x 0.978 where his UNIQUAC has none -- 0.044 K deep, and decisive. acetone06: so C1 caps at 97.3 mol \% acetone against his 99.9 \% spec. acetone07: so C2's recycle caps at x\_IPA 0.580 against his 0.650, and 11.6 \% of the isopropanol entering it leaves in the water product.

Both column ceilings were WRITTEN DOWN BEFORE the columns existed and are pinned by check\_azeotrope\_ceiling. This case is where they stop being two separate findings and become one plant.

DEPARTURES FROM HIS FIGURE 1, each argued rather than absorbed

(a) NO HYDROGEN IN THE COLUMNS. His C1 has a partial condenser and a vent (0.0465 kmol/h, carrying 0.0130 kmol/h H2 and 0.0335 kmol/h acetone); Choupo's column has a total condenser and no vapour distillate. A permanent gas fed to a total condenser is asked to leave as a liquid, and the engine rightly refuses (hydrogen has no liquid heat capacity at 330 K). So the separator's LIQUID goes to C1 and the hydrogen leaves with the offgas, i.e. the columns are modelled downstream of his vent. Cost, stated: the 0.0335 kmol/h of acetone his vent loses is not lost here, which flatters this plant's acetone recovery by 0.10 \%.

(b) THE SEPARATOR DOES THE COOLING. An isothermal flash at 318 K computes the duty to bring its feed there, so this one unit stands for his HX1 plus his separator and its Q is directly comparable to his 0.8993 MW. A separate `heater` was tried in acetone03 and refused, correctly: Choupo's heater takes a DUTY and reports the temperature, because a heater with a specified outlet temperature is a design specification.

(c) THE SEPARATOR PRESSURE IS ARGUED, NOT FITTED. Luyben states 1.5 atm for the absorber gas, which is downstream; there is no compressor between reactor and separator and pressure only falls, so the separator sits at the reactor's 2.6 atm. acetone03 records the sweep that would have tempted a fitted value and why it was refused.

(d) THE ABSORBER'S SECOND OUTLET REJOINS C1's FEED, as in his Figure 1 (his F1 of 72.85 kmol/h is the flash liquid AFTER the absorber bottoms of 20.05 kmol/h rejoin it).

(e) NO VENT/PURGE ON THE RECYCLE LOOP. Luyben has none either -- the loop closes through C2's overhead and nothing else -- so this is his structure, not a simplification. It does mean any component that both enters and cannot leave would accumulate without bound; here water leaves in three places and isopropanol in two, so the loop is closed and bounded.

\end{tutnarrative}
\tutreadme{These are the dictionaries of the closed acetone flowsheet described in [../acetone-plant-closure-state.md](../acetone-plant-closure-state.md). It does not converge yet, and the reason is one named engine gap (F3 in that record: no unit can remove a component completely).}
\tutreadme{tutorials/ is the verification and regression corpus. Everything in it is run by bin/runTests and is expected to answer. Putting a case there that cannot answer would have meant claiming one of two markers, and neither is true of this case:}
\tutreadme{$\ast$ .expect-nonconvergence --- every one of the five cases carrying it is a DELIBERATE TEACHING REFUSAL (flash15\_refused\_salt\_basis\_rank, tsa02\_refused\_capacity, ?): the refusal $\ast$is$\ast$ the lesson, and the case is finished. This one is not refusing anything on purpose; it is unfinished. $\ast$ .known-broken --- the corpus carries ZERO of these, and that zero is a claim the project makes about its own health. Spending it to park work in progress would be paying a real claim for my convenience.}
\tutreadme{So the case sits beside the record that explains it, makes no corpus claim, and moves to tutorials/plant/acetonePlant/ the day it converges --- at which point it gets a golden, a seal and an entry in the case manifest like any other case.}
\tutreadme{0/ is NOT kept: it is generated by bin/choupo-init0 from the authored inlets plus the tear seed, and a stale copy of a generated state is a second home for numbers the tool derives. The tear seed itself --- the one part of 0/ that is authored rather than propagated --- is described in the closure record.}
\tutreadme{The absorber's second thermodynamic world is declared inline in system/flowsheetDict as a per-unit thermo \{\} block, not as a separate file, because that is where the model boundary belongs and where the run announces it.}
\begin{tutgolden}
\goldrow{stream}{B1}{F}{0.0118946232842}
\goldrow{stream}{F1}{F}{0.0195617446099}
\goldrow{stream}{Rin}{F}{0.0160666666667}
\goldrow{stream}{Rout}{F}{0.0255036799795}
\goldrow{stream}{absBottoms}{F}{0.00608647733675}
\goldrow{stream}{acetoneProduct}{F}{0.00762772388347}
\goldrow{stream}{columnFeed}{F}{0.0195223471677}
\goldrow{stream}{flashGas}{F}{0.0120284127064}
\goldrow{stream}{flashLiquid}{F}{0.0134752672731}
\goldrow{stream}{freshFeed}{F}{0.0144333333333}
\goldrow{stream}{ipaRecycle}{F}{0.00163333333333}
\goldrow{stream}{mixedFeed}{F}{0.0160666666667}
\goldrow{stream}{offgas}{F}{0.0114974909252}
\goldrow{stream}{ventGas}{F}{3.9397442245e-05}
\goldmore{139}
\end{tutgolden}

\clearpage
\tutentry{esterification2sector}{tutorials/plant/esterification2sector}
\tutwhat{Esterification plant -- two sectors (REACTION + SEPARATION) sharing components; per-sector constant/ data (S1 kinetics, S2 NRTL pair); ethanol-water pair inherited from the standard library}
\tutreadme{A two-sector plant that exists to prove one thing: data lives with the sector that owns it, and a whole-plant run finds it there.}
\tutreadme{Components are shared across the plant. The kinetics belong to REACTION and are never copied to the root. The ethylAcetate-water NRTL pair belongs to SEPARATION --- a $\ast$particular$\ast$ datum, fitted for this separation --- while ethanol-water walks up to the standard catalogue. Three levels of ownership, one run.}
\tutreadme{That perNode is the whole case. The sector's pair beat the standard library, and the log says so rather than leaving you to wonder which one it used.}
\tutreadme{Feed 100 kmol/h equimolar acetic acid + ethanol at 360 K. The CSTR (V\_R = 0.5 m$^3$) reaches 34.2 \% conversion --- a PLACEHOLDER number: the kinetics in sectors/REACTION/constant/reactions are still the pseudo-first-order stand-in, not the 2nd-order fit this case's own property evidence supports (see CLAUDE.md, "Pending"). The flash at 355 K, 1.013 bar splits the product into a vapour rich in ethyl acetate (bp 350 K) and ethanol, and a liquid rich in water and the unconverted acid. Mass closes exactly: 34.18 + 65.82 = 100.00 kmol/h.}
\tutreadme{This case shipped for weeks marked .known-broken, and it earned the mark twice over. Its flash carried P 1.01325; with no unit, which is 1.01 pascal --- a hard vacuum --- so the K-values exploded and everything went to the vapour. And it was never wired: the sectors existed as thermo regions with the connections commented out, awaiting a "curation phase" that never ended.}
\tutreadme{Fixing it surfaced two real engine bugs, both of the same shape --- data owned by a $\ast$sector$\ast$ was invisible from a whole-plant run:}
\begin{tutgolden}
\goldrow{stream}{Feed}{F}{0.0277777777778}
\goldrow{stream}{liquid}{F}{0.0114069310496}
\goldrow{stream}{reactorOut}{F}{0.0277777777778}
\goldrow{stream}{vapor}{F}{0.0163708467282}
\goldrow{utility}{SEPARATION.flash}{cooling.coolingWater.duty\_kW}{-229.173881338}
\goldrow{utility}{SEPARATION.flash}{cooling.coolingWater.T}{355}
\goldrow{utility}{SEPARATION.flash}{cooling.coolingWater.kg\_s}{5.48262874014}
\goldrow{utility}{SEPARATION.flash}{cooling.coolingWater.eur\_h}{0.330010389127}
\end{tutgolden}

\clearpage
\tutentry{hda}{tutorials/plant/hda}
\begin{tutscheme}
\begin{tikzpicture}[>={Stealth[length=2mm]}, every node/.style={font=\scriptsize}]
  \node[draw, rounded corners=2pt, align=center, minimum width=1.9cm, minimum height=0.85cm, fill=black!4] (Mixer) at (60.9,-0.0) {\textbf{Mixer}\\[-1pt]{\tiny mixer}};
  \node[draw, rounded corners=2pt, align=center, minimum width=1.9cm, minimum height=0.85cm, fill=black!4] (Reactor) at (63.8,-0.0) {\textbf{Reactor}\\[-1pt]{\tiny conversionReactor}};
  \node[draw, rounded corners=2pt, align=center, minimum width=1.9cm, minimum height=0.85cm, fill=black!4] (HPFlash) at (66.7,-0.0) {\textbf{HPFlash}\\[-1pt]{\tiny isothermalFlash}};
  \node[draw, rounded corners=2pt, align=center, minimum width=1.9cm, minimum height=0.85cm, fill=black!4] (GasSplit) at (69.6,-0.0) {\textbf{GasSplit}\\[-1pt]{\tiny splitter}};
  \node[draw, rounded corners=2pt, align=center, minimum width=1.9cm, minimum height=0.85cm, fill=black!4] (Column) at (69.6,-1.7) {\textbf{Column}\\[-1pt]{\tiny shortcutColumn}};
  \draw[->] ($(Mixer.west)+(-1.05,-0.17)$) -- ($(Mixer.west)+(0,-0.17)$) node[at start, below, font=\tiny, anchor=north west] {TolueneFeed};
  \draw[->] ($(Mixer.west)+(-1.05,0.17)$) -- ($(Mixer.west)+(0,0.17)$) node[at start, above, font=\tiny, anchor=south west] {H2Makeup};
  \draw[->, dashed] (GasSplit.south) .. controls +(0,-1.0) and +(0,-1.0) .. (Mixer.south) node[midway, below, font=\tiny] {H2Recycle};
  \draw[->, dashed] (Column.south) .. controls +(0,-1.0) and +(0,-1.0) .. (Mixer.south) node[midway, below, font=\tiny] {TolueneRecycle};
  \draw[->] (Mixer) -- (Reactor) node[midway, above, font=\tiny] {mixed};
  \draw[->] (Reactor) -- (HPFlash) node[midway, above, font=\tiny] {Effluent};
  \draw[->] (HPFlash) -- (GasSplit) node[midway, above, font=\tiny] {hpVapor};
  \draw[->] (HPFlash) -- (Column) node[midway, above, font=\tiny] {hpLiquid};
  \draw[->] ($(GasSplit.east)+(0,0.00)$) -- ($(GasSplit.east)+(1.05,0.00)$) node[at end, above, font=\tiny, anchor=south east] {Purge};
  \draw[->] ($(Column.east)+(0,0.00)$) -- ($(Column.east)+(1.05,0.00)$) node[at end, above, font=\tiny, anchor=south east] {Benzene};
\end{tikzpicture}
\end{tutscheme}
\tutwhat{HDA plant -- toluene + H2 $\to$ benzene + CH4. A single FLAT flowsheet with two plant-level recycles (an H2 gas loop and an unreacted-toluene liquid loop); cross-sector tears are unsupported, so it is not fractal. Furnace lumped into the isothermal reactor; light-gas K-values rest on Antoine extrapolation (see thermoPhysPropDict).}
\begin{tutgolden}
\goldrow{stream}{Benzene}{F}{0.0308570333668}
\goldrow{stream}{Effluent}{F}{0.589422462034}
\goldrow{stream}{H2Makeup}{F}{0.0305555555556}
\goldrow{stream}{H2Recycle}{F}{0.522049697299}
\goldrow{stream}{Purge}{F}{0.0274762998579}
\goldrow{stream}{TolueneFeed}{F}{0.0277777777778}
\goldrow{stream}{TolueneRecycle}{F}{0.00903943150968}
\goldrow{stream}{hpLiquid}{F}{0.0398964648764}
\goldrow{stream}{hpVapor}{F}{0.549525997157}
\goldrow{stream}{mixed}{F}{0.589422462034}
\goldrow{kpi}{Column}{B}{0.00903943150968}
\goldrow{kpi}{Column}{D}{0.0308570333668}
\goldrow{kpi}{Column}{N\_min}{8.3953506183}
\goldrow{kpi}{Column}{N\_theoretical}{20.1717591626}
\goldmore{38}
\end{tutgolden}

\clearpage
\tutentry{lithiumBrinePlant}{tutorials/plant/lithiumBrinePlant}
\tutwhat{lithiumBrinePlant -- THE reference case: a FRACTAL lithium-carbonate plant from salar brine, WIRED into one solver. Five SECTORS, each a thermo WORLD of its own -- BRINE (electrolyte/Pitzer), EXTRACTION (gamma-gamma NRTL LLE), CARBONATION (Li2CO3 retrograde solubility), HOTAIR (Gibbs combustion), FINISHING (drying + cyclone) -- connected sector to sector, each sector declaring its own propertyPackage (its world) on the global component set, with optional per-unit thermo overrides where a single unit needs a different world. The proof that the architecture is fractal: sectors are composite boxes, streams flow between them, model boundaries audited.}
\begin{tutnarrative}
Five REAL sectors. Streams are named physical graph edges. Sector names are subdomains; unit operations live inside each sector.

\end{tutnarrative}
\begin{tutgolden}
\goldrow{stream}{carbLiquor}{F}{0.0277777777778}
\goldrow{stream}{cleanAir}{F}{0.0320991487231}
\goldrow{stream}{emulsion}{F}{0.0380203541089}
\goldrow{stream}{fines}{F}{0}
\goldrow{stream}{fuelAir}{F}{0.0108555555558}
\goldrow{stream}{halite}{F}{0.000868534779985}
\goldrow{stream}{hotAir}{F}{0.010855555556}
\goldrow{stream}{humidExhaust}{F}{0.0320991487231}
\goldrow{stream}{liquor}{F}{0.0269092429978}
\goldrow{stream}{loadedOrganic}{F}{0.0124069292995}
\goldrow{stream}{organic}{F}{0.0111111111111}
\goldrow{stream}{product}{F}{0.00653418461064}
\goldrow{stream}{raffinate}{F}{0.0256134248094}
\goldrow{stream}{salarBrine}{F}{0.0277777777778}
\goldmore{1}
\end{tutgolden}

\subsection{polycaprolactonePlant}
\grouptheory{polycaprolactonePlant}
\clearpage
\tutentry{properties}{tutorials/plant/polycaprolactonePlant/properties}
\begin{tutscheme}
\begin{tikzpicture}[>={Stealth[length=2mm]}, every node/.style={font=\scriptsize}]
  \node[draw, rounded corners=2pt, align=center, minimum width=3.1cm, minimum height=0.8cm, fill=black!4] at (0.0,-0.0) {\textbf{pclRepeat\_Tg\_yang2020}\\[-1pt]{\tiny estimateComponent}};
\end{tikzpicture}
\end{tutscheme}
\tutwhat{Close the loop: ESTIMATE the PCL repeat-unit glass transition Tg by the Yang 2020 group-contribution scheme --- the plant gives Mn/PDI, this gives the physical property from the SAME repeat unit -[O-(CH2)5-CO]-}
\begin{tutnarrative}
pclProperties -- CLOSE THE LOOP: process $\to$ molar mass $\to$ physical property.

The plant (../) gave you the chain statistics from the KINETICS: Mn, Mw, PDI on the polycaprolactone it makes. Here you ESTIMATE the polymer's PHYSICAL PROPERTIES from the SAME repeat unit -[O-(CH2)5-CO]- by group contribution --- the bridge from "how big are the chains" to "how does the material behave".

TWO estimators, two properties, on the repeat unit pclRepeat (M0 = 114.14):

(A) Tg, the GLASS-TRANSITION temperature --- Yang 2020 (run below). Repeat unit -[O-(CH2)5-CO]- decomposes into Yang main-chain groups: O x1 : MW 16 Yg -14.718 CH2 x5 : MW 14 Yg 4.026 (each) CO x1 : MW 28 Yg 4.370

M0 = 16 + 5$\ast$14 + 28 = 114 g/mol (matches the reactor M0!) Yg(inf) = -14.718 + 5$\ast$4.026 + 4.370 = 9.782 (1e3 g.K/mol) Tg(inf) = 9.782e3 / 114 \textasciitilde{} 86 K (the additive estimate) (Measured PCL Tg \textasciitilde{} 213 K / -60 C; an additive main-chain scheme is weak for a flexible aliphatic polyester --- the model tells the truth about itself, which IS the lesson. See the console's per-cent error vs the anchor.)

(B) DENSITY --- Van Krevelen (Bondi Vw). NOT runnable on PCL with the shipped data: the Slice-1 Van Krevelen set ships only hydrocarbon + chloro groups (CH3 CH2 CH C ACH AC CHCl CH2Cl); it has NO ester -O- / -CO- group, so the repeat unit -[O-(CH2)5-CO]- cannot be decomposed. Running it would HARD- ERROR with "unknown polymer group 'O'" --- the estimator refuses to invent a value (no silent zero). The honest path: density needs the ester groups curated into data/standards/parameters/vanKrevelen.dat first. The worked Van Krevelen DENSITY example lives in tutorials/props/estimate/vanKrevelen01\_polystyrene\_density (a pure- hydrocarbon repeat unit the Slice-1 set CAN do).

This is the same M0 = 114.14 the reactor's polymer\{\} block used for Mn = M0/(1-p) --- one repeat-unit mass, two places (process AND property).

\end{tutnarrative}
\tutreadme{A choupoProps sub-case of [../](..) (the polycaprolactone plant). The plant gives you the chain statistics from the kinetics (Mn / Mw / PDI); here you estimate the polymer's physical property --- the glass transition Tg --- from the same repeat unit -[O-(CH2)5-CO]- (M0 = 114.14, the exact number the reactor used for Mn = M0/(1-p)).}
\tutreadme{$\ast$ Tg via Yang 2020 (CC-BY) --- runnable. Tg(inf) = 85.8 K; the console prints the -59.7 \% deviation from the measured \textasciitilde{}213 K (an honest limit of an additive scheme on a flexible aliphatic polyester). M0 = 114 is goldened (an exact mass-balance fact); Tg is shown but not goldened as validated physics. $\ast$ Density via Van Krevelen --- $\ast$not$\ast$ runnable on PCL with the shipped data (the Slice-1 group set has no ester -O-/-CO- group, so it would hard-error rather than invent a value). See the full explanation and the runnable hydrocarbon example linked from the parent [../README.md](../README.md).}
\tutreadme{The full teaching narrative lives in the parent README. See [system/propsDict](system/propsDict) for the exact group decomposition and the estimator invocation.}
\begin{tutgolden}
\goldrow{diag}{pclRepeat\_Tg\_yang2020}{M0\_g\_per\_mol}{114}
\goldrow{diag}{pclRepeat\_Tg\_yang2020}{Tg\_K}{85.81}
\goldrow{diag}{pclRepeat\_Tg\_yang2020}{YgSum\_1e3gKmol}{9.782}
\end{tutgolden}

\clearpage
\tutentry{sugarPlantEconomicsSweep}{tutorials/plant/sugarPlantEconomicsSweep}
\tutwhat{Sugar plant ECONOMICS SWEEP -- IRR / payback / NPV vs RawJuice sucrose content (Phase-2 SweepDriver post-chain differentiator)}
\begin{tutnarrative}
/ system/flowsheetDict (PLANT level -- COMPOSITE)

The flagship sugar plant (the ChemicalPlantTutorial flowsheet, reused verbatim -- sectors and units included), here driven by system/outerDict: a SweepDriver walks the RawJuice sucrose fraction and runs the FULL post chain (sizing $\to$ costing $\to$ economics, system/postDict) on every point, so IRR / NPV / payback vs feed quality come out as sweep responses.

A composite node: it has `children` (sub-flowsheets, here the sectors) and `connections` (the pipework between them). The same shape a sector uses to list its units -- one level up. This is the fractal.

Naming convention (see docs/ai/case-layout.md): SECTORS in CAPS : CONCENTRATION, DRYING, FERMENTATION Units in PascalCase : Cryst, Evap1, Evap2 Streams in PascalCase : RawJuice, Magma, Powder, EvapVapour Components in lowercase : water, sucrose, N2

\end{tutnarrative}
\tutreadme{This is the flagship sugar plant (tutorials/plant/ChemicalPlantTutorial) driven by a SweepDriver instead of a single pass. It exists to show the one thing a single CAPEX number cannot: how the project economics move when a process variable changes.}
\tutreadme{The driver runs the FULL post-processing chain (sizing $\to$ costing $\to$ economics, in system/postDict) on every converged point, so economics.IRR, economics.paybackYears, economics.NPV, ? resolve as ordinary KPI look-ups (the Phase-2 SweepDriver post-chain wiring). The result lands in sweep\_results.csv.}
\tutreadme{Sweeping the RawJuice sucrose content 0.09 $\to$ 0.16 mol fraction walks the plant from sub-economic, through break-even, to clearly profitable:}
\tutreadme{Richer juice $\to$ more crystalline sugar per unit feed $\to$ more revenue against an almost-unchanged plant cost $\to$ the IRR climbs and the payback appears and shortens. The NPV crosses zero between 0.104 and 0.118.}
\tutreadme{Both use the same constant/economics prices (EU Sugar Market Observatory white sugar, ECB EUR/USD, $\ast$Chemical Engineering$\ast$ CEPCI, US statutory tax) and the same system/postDict cost chain --- only the run mode differs.}

\clearpage
\tutentry{tartaricAcid}{tutorials/plant/tartaricAcid}
\begin{tutscheme}
\begin{tikzpicture}[>={Stealth[length=2mm]}, every node/.style={font=\scriptsize}]
  \node[draw, rounded corners=2pt, align=center, minimum width=1.9cm, minimum height=0.85cm, fill=black!4] (M1) at (0.0,-0.0) {\textbf{M1}\\[-1pt]{\tiny mixer}};
  \node[draw, rounded corners=2pt, align=center, minimum width=1.9cm, minimum height=0.85cm, fill=black!4] (R1) at (2.9,-0.0) {\textbf{R1}\\[-1pt]{\tiny conversionReactor}};
  \node[draw, rounded corners=2pt, align=center, minimum width=1.9cm, minimum height=0.85cm, fill=black!4] (M2) at (5.8,-0.0) {\textbf{M2}\\[-1pt]{\tiny mixer}};
  \node[draw, rounded corners=2pt, align=center, minimum width=1.9cm, minimum height=0.85cm, fill=black!4] (R2) at (8.7,-0.0) {\textbf{R2}\\[-1pt]{\tiny conversionReactor}};
  \draw[->] ($(M1.west)+(-1.05,-0.17)$) -- ($(M1.west)+(0,-0.17)$) node[at start, below, font=\tiny, anchor=north west] {bitartrate};
  \draw[->] ($(M1.west)+(-1.05,0.17)$) -- ($(M1.west)+(0,0.17)$) node[at start, above, font=\tiny, anchor=south west] {lime};
  \draw[->] (M1) -- (R1) node[midway, above, font=\tiny] {slurry};
  \draw[->] (R1) -- (M2) node[midway, above, font=\tiny] {precipOut};
  \draw[->] ($(M2.west)+(-1.05,0.00)$) -- ($(M2.west)+(0,0.00)$) node[at start, above, font=\tiny, anchor=south west] {sulfuric};
  \draw[->] (M2) -- (R2) node[midway, above, font=\tiny] {acidFeed};
  \draw[->] ($(R2.east)+(0,0.00)$) -- ($(R2.east)+(1.05,0.00)$) node[at end, above, font=\tiny, anchor=south east] {product};
\end{tikzpicture}
\end{tutscheme}
\tutwhat{Tartaric-acid plant (wine route): lime precipitation of calcium tartrate, then sulfuric acidulation to tartaric acid + gypsum. Molecular reactors; the solution chemistry sub-study is tutorials/props/electrolyte/tartaricAcid\_acidulation. HONEST LIMITS: H2SO4 declared nonvolatile by case overlay (Psat \textasciitilde{}1e-5 Pa at process T); precipitates ride as nonvolatile solutes with NO rigorous solid-liquid separation; reactor duties honestly UNAVAILABLE (species without ideal-gas Cp -- a named gap, never a fake zero).}
\begin{tutgolden}
\goldrow{stream}{acidFeed}{F}{0.943266666667}
\goldrow{stream}{bitartrate}{F}{0.583333333333}
\goldrow{stream}{lime}{F}{0.166666666667}
\goldrow{stream}{precipOut}{F}{0.7766}
\goldrow{stream}{product}{F}{0.943266666667}
\goldrow{stream}{slurry}{F}{0.75}
\goldrow{stream}{sulfuric}{F}{0.166666666667}
\goldrow{kpi}{M1}{F\_out}{0.75}
\goldrow{kpi}{M1}{P\_out}{100000}
\goldrow{kpi}{M1}{T\_out}{333}
\goldrow{kpi}{M1}{n\_in}{2}
\goldrow{kpi}{M1}{vf}{0}
\goldrow{kpi}{M2}{F\_out}{0.943266666667}
\goldrow{kpi}{M2}{P\_out}{100000}
\goldmore{13}
\end{tutgolden}

\subsection{twoSectorDemo}
\grouptheory{twoSectorDemo}
\clearpage
\tutentry{SECTOR\_R}{tutorials/plant/twoSectorDemo/SECTOR\_R}
\begin{tutscheme}
\begin{tikzpicture}[>={Stealth[length=2mm]}, every node/.style={font=\scriptsize}]
  \node[draw, rounded corners=2pt, align=center, minimum width=1.9cm, minimum height=0.85cm, fill=black!4] (R1) at (0.0,-0.0) {\textbf{R1}\\[-1pt]{\tiny stoichReactor}};
  \draw[->] ($(R1.west)+(-1.05,0.00)$) -- ($(R1.west)+(0,0.00)$) node[at start, above, font=\tiny, anchor=south west] {feed};
  \draw[->] ($(R1.east)+(0,0.00)$) -- ($(R1.east)+(1.05,0.00)$) node[at end, above, font=\tiny, anchor=south east] {rOut};
\end{tikzpicture}
\end{tutscheme}
\tutwhat{SECTOR\_R -- the reactor sector. Runs its OWN compiled unit op (stoichReactor) on the plant-level thermo (cascaded from ../constant).}
\begin{tutnarrative}
one custom reactor (compiled from code/). Build + run with: bin/buildCode tutorials/plant/twoSectorDemo/SECTOR\_R runCase tutorials/plant/twoSectorDemo/SECTOR\_R ---

\end{tutnarrative}

\clearpage
\tutentry{SECTOR\_S}{tutorials/plant/twoSectorDemo/SECTOR\_S}
\begin{tutscheme}
\begin{tikzpicture}[>={Stealth[length=2mm]}, every node/.style={font=\scriptsize}]
  \node[draw, rounded corners=2pt, align=center, minimum width=1.9cm, minimum height=0.85cm, fill=black!4] (C1) at (0.0,-0.0) {\textbf{C1}\\[-1pt]{\tiny sharpSplitColumn}};
  \draw[->] ($(C1.west)+(-1.05,0.00)$) -- ($(C1.west)+(0,0.00)$) node[at start, above, font=\tiny, anchor=south west] {feed};
  \draw[->] ($(C1.east)+(0,-0.17)$) -- ($(C1.east)+(1.05,-0.17)$) node[at end, below, font=\tiny, anchor=north east] {top};
  \draw[->] ($(C1.east)+(0,0.17)$) -- ($(C1.east)+(1.05,0.17)$) node[at end, above, font=\tiny, anchor=south east] {bottom};
\end{tikzpicture}
\end{tutscheme}
\tutwhat{SECTOR\_S -- the separation sector. Runs its OWN compiled unit op (sharpSplitColumn) on the plant-level thermo (cascaded from ../constant).}
\begin{tutnarrative}
one custom distillation (compiled from code/). Build + run: bin/buildCode tutorials/plant/twoSectorDemo/SECTOR\_S runCase tutorials/plant/twoSectorDemo/SECTOR\_S ---

\end{tutnarrative}

